% PROPOZYCJA do ręcznego włączenia po weryfikacji promotorskiej.
% Zachowuje cztery hipotezy z rozdziału 1 PDF i porządkuje tabelę 5.2.
% Propozycja do ręcznego włączenia po weryfikacji promotorskiej.
% Oceny są wstępne i ograniczone do testowanego laboratorium; przed finalizacją
% zweryfikować mapowanie na pytania badawcze i całość artefaktów.

\subsection{Proponowane ujednolicenie hipotez badawczych}
\label{subsec:hipotezy-propozycja}

W celu zachowania spójności metodyki przyjęto cztery hipotezy przedstawione
w rozdziale pierwszym. Zagadnienia segmentacji sieci, wykorzystania danych-canary
oraz transportu MCP są analizowane jako warunki lub scenariusze badawcze, a nie
jako dodatkowe hipotezy.

\begin{table}[ht]
\centering
\small
\caption{Proponowane kryteria weryfikacji hipotez badawczych}
\label{tab:hipotezy-weryfikacja-propozycja}
\begin{tabular}{|p{0.18\textwidth}|p{0.25\textwidth}|p{0.31\textwidth}|p{0.16\textwidth}|}
\hline
\textbf{Hipoteza} & \textbf{Kryterium operacyjne} &
\textbf{Planowane źródła dowodowe} & \textbf{Ocena po badaniach} \\
\hline
H1: Telemetria pochodząca wyłącznie z agentowego środowiska programistycznego
nie pozwala na pełne potwierdzenie rzeczywistych skutków operacji agenta. &
Porównanie informacji dostępnych w Kilo z niezależnym potwierdzeniem wykonania
i skutku. Wskazanie etapów łańcucha, których sam zapis aplikacyjny nie potwierdza. &
Kilo, Gateway, MCP, auditd, Target i Zeek/PCAP; w SC-03 rep01 UI pokazało
odmowę, a Gateway zarejestrował \texttt{policy\_load\_error}, co wymagało korekty
interpretacji. SC-05/SC-06-A zestawiono z wynikiem Target. &
Proponowana ocena: częściowo potwierdzona; widok Kilo sam nie wystarcza do
ustalenia przyczyny i niezależnego potwierdzenia skutku. \\
\hline
H2: Wykonanie tego samego niepożądanego działania w szerokim kontekście
użytkownika może powodować większy skutek niż wykonanie go w wydzielonym
kontekście o ograniczonych uprawnieniach. &
Porównanie jednego odczytu syntetycznego chronionego zasobu przez
\texttt{agent-wide} i \texttt{agent-limited} przy niezmienionym zadaniu, konfiguracji MCP
i pliku. Oczekiwana para: odczyt dozwolony i odmowa systemu operacyjnego. &
Trzy niezależne pary Kilo: \texttt{H2-KILO-REP01-20261010},
\texttt{H2-KILO-REP02-20261010} i \texttt{H2-KILO-REP03-20261010}. Wszystkie: wide
success, limited permission denied; MCP JSONL, auditd (openat success/no,
limited exit=-13) i odbiór na Monitor. &
Proponowana ocena: częściowo potwierdzona w zakresie badanego pliku i POSIX ACL. \\
\hline
H3: Deterministyczna kontrola narzędzi i parametrów ich wywołań ogranicza
liczbę skutecznie wykonanych operacji niedozwolonych, również w przypadku
manipulacji kontekstem przetwarzanym przez model. &
SC-02-A allow 3/3 z odbiorem; SC-02-B ręczne deny 3/3 bez attempt/odbioru;
SC-03 policy deny 3/3 bez attempt/odbioru; SC-06-A dodatnia kontrola 3/3.
SC-06-D w 3/3 nie wygenerowało żądania submit, więc kontroler w tym wariancie
nie został wywołany. &
Gateway decision/approval/attempt/result, Target receiver, PCAP/Zeek dla
pozytywnych kontroli. &
Proponowana ocena: częściowo potwierdzona; kontrola działała dla jawnych
żądań, a warunek pośredniej manipulacji nie uruchomił jej. \\
\hline
H4: Korelacja telemetrii aplikacyjnej, hosta i sieciowej umożliwia pełniejszą
rekonstrukcję działania agenta niż wykorzystanie dowolnego z tych źródeł osobno. &
Porównanie etapów rekonstrukcji dostępnych w każdym źródle oddzielnie
i po ich zestawieniu; jawne oznaczenie brakujących identyfikatorów i poziomu
pewności korelacji. &
SC-01/SC-04: lokalny MCP, auditd i Monitor; SC-05: Gateway, Target journal,
PCAP/Zeek; SC-06-A: Gateway, receiver i PCAP/Zeek. Część korelacji Target/PCAP
opiera się na czasie i adresach, bo run ID nie jest tam propagowany. &
Proponowana ocena: częściowo potwierdzona; źródła uzupełniają obraz, lecz
atrybucja nie zawsze jest bezpośrednia. \\
\hline
\end{tabular}
\end{table}

Wynik pojedynczej próby nie jest samodzielnie oceną hipotezy. W tabeli należy
wskazać przebiegi stanowiące podstawę oceny, liczbę prób w każdym warunku,
odchylenia od protokołu oraz ograniczenia wnioskowania. Pilot ACL
\texttt{KILO-H2-KILO-20261010-01} nie wchodzi do licznika trzech par.
Trzy pary dały ten sam wynik dla badanego syntetycznego pliku i dwóch kont,
co wspiera H2 w tym ograniczonym warunku. Nie dowodzi to ogólnej izolacji
ani skuteczności kontroli alternatywnych dróg dostępu; profile Kilo zawierały
\texttt{bash: allow}, a seria wymuszała bezpośrednio jedno wywołanie MCP.

% Koniec propozycji. Przed kompilacją sprawdzić, czy preambuła dokumentu
% obsługuje kodowanie UTF-8 oraz czy szerokość tabeli mieści się w szablonie.
